\section{Introduction}

GPU-heavy serving systems, particularly large language model~(LLM) inference engines, can
experience performance problems for many reasons: request queuing, batching delays, GPU
compute saturation, GPU memory pressure, long input or output sequences, or transient
workload shifts.
Operators typically have access to cheap, always-on telemetry---request latency, throughput,
GPU utilization, and GPU memory utilization---that can reveal \emph{that} something is wrong.
These cheap signals, however, frequently cannot explain \emph{why} a system is slow: a spike
in p95 latency can be produced by queueing delay, by a compute-bound kernel, by memory-bound
data movement, or by some combination of the three, and cheap utilization counters alone
cannot distinguish between these causes.

GPU profiling tools such as NVIDIA Nsight Systems can resolve this ambiguity by exposing
kernel-level evidence: execution timing, launch frequency, and duration distributions for
individual GPU kernels.
This evidence is detailed enough to identify the true bottleneck class, but heavy profiling
is also expensive to collect and disruptive to run.
In current practice, GPU profiling is almost always a manual activity: an operator notices a
problem, starts a profiler with a fixed-duration window, waits for the window to elapse,
and then inspects the resulting trace.
This workflow has four practical limitations.
First, it requires a human in the loop, which does not scale to automated monitoring
pipelines.
Second, it relies on a fixed profiling window that is set in advance and cannot adapt to how
quickly the evidence actually converges.
Third, a fixed window can easily collect more evidence than necessary, wasting profiler
overhead and producing oversized trace artifacts.
Fourth, and most fundamentally, manual fixed-window profiling has no principled way to
answer when enough evidence has been collected---it stops because the timer expired, not
because the diagnosis is complete.

This last limitation motivates the central problem studied in this paper:

\begin{quote}
\emph{When has a GPU tracing system collected sufficient evidence to make a diagnosis, and
when is additional heavy tracing no longer worth its cost?}
\end{quote}

We refer to this as the problem of \emph{diagnostic sufficiency}.
It is distinct from, and complementary to, the problem of trace volume reduction studied by
prior work on adaptive distributed tracing.
Adaptive tracing systems for distributed request-level traces decide which requests or spans
to sample, but they generally operate above the GPU and do not reason about GPU-specific
evidence tiers such as utilization counters, memory pressure, or kernel-level timing.
Conversely, GPU profiling tools expose rich low-level evidence but are built as manual,
fixed-window instruments; they do not decide when profiling should start, continue, or stop.
The gap between these two bodies of work is a missing runtime control problem: managing the
lifecycle of heavy GPU tracing automatically, from activation through evidence sufficiency to
de-escalation.

\subsection{Approach}

This paper proposes and evaluates an \emph{adaptive GPU tracing controller} that closes this
gap.
The controller organizes evidence into two cost tiers.
Cheap evidence---request latency, throughput, queueing delay, GPU utilization, and GPU memory
utilization---is collected continuously at negligible cost.
Heavy evidence---kernel execution timing collected via short Nsight Systems profiler
bursts---is collected only when the cheap evidence is insufficient to explain an observed
anomaly.
After each burst, the controller evaluates whether the resulting kernel-level diagnosis has
stabilized.
If additional bursts would no longer change the diagnosis, or if cheap signals show sustained
recovery, the controller de-escalates and returns to cheap-only monitoring.
We design and compare six concrete short-burst stopping policies that instantiate this idea
with different cost--reliability tradeoffs, ranging from a single fixed burst to a hybrid
policy that requires both diagnosis stability and workload recovery before stopping.

We implement and evaluate this controller on a realistic GPU-heavy serving system: vLLM
serving the \texttt{Qwen2.5-7B-Instruct} model on NVIDIA L4 and A100 GPUs, across six
distinct serving scenarios (healthy operation, queue pressure, long prompts, long outputs,
compute saturation, and KV-cache pressure) and three controlled CUDA/PyTorch microbenchmarks
with known ground-truth bottleneck labels.

\subsection{Research Questions}

We structure our evaluation around five research questions that together examine whether,
how accurately, at what cost, and through which mechanism an adaptive controller can manage
GPU trace de-escalation.

\paragraph{RQ1: Can an adaptive GPU tracing controller determine when heavy GPU tracing has
collected sufficient diagnostic evidence?}
This is the central research question of the paper. It asks whether the controller can stop
heavy tracing automatically, once additional kernel-timing evidence no longer changes the
diagnosis, while still producing a correct diagnosis. We evaluate whether the controller
stops earlier than a fixed manual profiling window, whether it avoids stopping prematurely,
and whether it preserves diagnosis quality while doing so.

\paragraph{RQ2: How accurate is automatic GPU tracing compared with a manual profiling
approach?}
Reducing tracing duration is only useful if it does not come at the cost of diagnostic
accuracy. RQ2 directly compares the automatic controller against a manual fixed-window
baseline using expected-label match rate, ambiguous-diagnosis rate, and time-to-diagnosis,
to establish whether automatic tracing is a safe substitute for manual profiling rather than
merely a cheaper one.

\paragraph{RQ3: Does automatic GPU tracing introduce overhead compared with manual profiling,
and is the overhead negligible?}
An adaptive controller must itself be cheap to run continuously. RQ3 measures two forms of
overhead: the profiler-cost overhead of heavy tracing relative to the fixed-window baseline,
and the runtime perturbation introduced by the controller's own always-on cheap-metric
collection, measured as request-latency and throughput impact on the serving workload.

\paragraph{RQ4: Which short-burst GPU tracing policies work best?}
Diagnostic sufficiency can be operationalized through different stopping policies, each
making a different cost--reliability tradeoff. RQ4 compares six candidate policies---Fixed
Burst, Repeated Fixed Burst, Stability Stop, Marginal Utility Stop, Counter-Recovery Stop,
and Hybrid Stop---on diagnosis accuracy, tracing duration, premature-stop rate, and
re-escalation rate, including under a controlled ambiguity stress condition that tests
robustness when the first burst of evidence is inconclusive.

\paragraph{RQ5: Which runtime signals are most useful for deciding when to stop heavy GPU
tracing?}
The policies in RQ4 are built from underlying runtime signals: diagnosis stability, latency
recovery, throughput recovery, GPU utilization recovery, and others. RQ5 isolates these
signals and evaluates each one's precision and reliability at identifying correct stopping
points, while also testing whether each signal can be fooled by ambiguous evidence. This RQ
is motivated by the observation that the best signal for \emph{starting} heavy tracing may
not be the best signal for \emph{stopping} it.

\subsection{Contributions}

This paper makes the following contributions.

\begin{enumerate}
  \item A formulation of automatic GPU trace de-escalation as a runtime control problem,
        distinct from existing adaptive distributed tracing and manual GPU profiling.
  \item A short-burst GPU tracing controller that activates and stops heavy kernel-level
        profiling without manual intervention, built on a two-tier evidence model.
  \item Six short-burst tracing policies, evaluated for diagnosis accuracy, tracing cost,
        premature-stop rate, and re-escalation rate, including under a controlled ambiguity
        stress condition.
  \item An evaluation against a manual fixed-window profiling baseline across controlled
        microbenchmarks and six realistic vLLM serving scenarios on two GPU platforms.
  \item Evidence that automatic tracing matches manual diagnosis accuracy while reducing
        profiler duration by 18.9\%--48.5\% and captured kernel instances by
        42.9\%--78.2\%, with negligible runtime overhead.
  \item A signal-level analysis identifying diagnosis stability as the most reliable stopping
        signal, and characterizing why several intuitive recovery-based signals are not safe
        to use as standalone stopping criteria.
\end{enumerate}

The remainder of this paper is organized as follows.
Section~\ref{sec:methodology} describes the controller design, evidence tiers, stopping
policies, and experimental implementation.
Section~\ref{sec:evaluation} presents the evaluation, organized by research question.
Later sections discuss related work, threats to validity, and conclusions.
